\chapter{RK4 + generic time integrator}
\label{chap:dev:rk4--generic-time-integrator}

\begin{quote}
Phase metadata. Spec dir:
\passthrough{\lstinline!specs/2026-05-02-phase-6-rk4-integrator/!}. Adds
\passthrough{\lstinline!solve::RK4!} integrator,
\passthrough{\lstinline!solve::Integrator.hpp!} aggregator,
tag-dispatched \passthrough{\lstinline!solve::integrate(...)!}.
Refactors \passthrough{\lstinline!solve::diffuse!} to template on the
integrator tag and \passthrough{\lstinline!solve::explicitEuler!} to be
a thin wrapper. Version target: \passthrough{\lstinline!0.6.0!}.
\end{quote}

\subsection{1. Theory}\label{theory}

The Runge--Kutta family generalizes forward Euler by evaluating the RHS
at carefully chosen intermediate points within each step and combining
the results to match more terms of the Taylor expansion of the exact
flow. Higher-order RK schemes have:

\begin{itemize}
\tightlist
\item
  \textbf{Higher accuracy} per step (local truncation error
  \passthrough{\lstinline!O(dt\^\{p+1\})!}, global error
  \passthrough{\lstinline!O(dt\^p)!}, where \passthrough{\lstinline!p!}
  is the method order).
\item
  \textbf{Larger stability regions} in the complex plane, so the method
  tolerates larger \passthrough{\lstinline!dt!} for
  stiff-but-not-too-stiff problems.
\item
  \textbf{Multiple RHS evaluations} per step (one per stage), making
  them more expensive per step.
\end{itemize}

The classic 4-stage Runge--Kutta method (RK4) has order
\passthrough{\lstinline!p = 4!} and uses 4 RHS evaluations per step. Its
Butcher tableau is

\[
\begin{array}{c|cccc}
0   &       &       &       &       \\
1/2 & 1/2   &       &       &       \\
1/2 & 0     & 1/2   &       &       \\
1   & 0     & 0     & 1     &       \\
\hline
    & 1/6   & 1/3   & 1/3   & 1/6
\end{array}
\]

Per-step update for \(\dot\psi = \mathrm{rhs}(\psi)\):

\[
\begin{aligned}
k_1 &= \mathrm{rhs}(\psi^n) \\
k_2 &= \mathrm{rhs}(\psi^n + (dt/2)\,k_1) \\
k_3 &= \mathrm{rhs}(\psi^n + (dt/2)\,k_2) \\
k_4 &= \mathrm{rhs}(\psi^n + dt\,k_3) \\
\psi^{n+1} &= \psi^n + (dt/6)(k_1 + 2 k_2 + 2 k_3 + k_4).
\end{aligned}
\]

For the heat equation \(\partial_t\psi = D\nabla^2\psi\), RK4 advances
each trig eigenfunction by the truncated Taylor series of
\(e^{D\,\lambda_h\,dt}\) --- exact through fourth order in the per-step
exponent.

\subsection{2. Math derivation}\label{math-derivation}

\subsubsection{RK4's order-4 accuracy}\label{rk4s-order-4-accuracy}

Apply RK4 to a linear scalar test problem \(\dot y = z\,y\) with
\(z = D\,\lambda_h\) (the discrete eigenvalue along one trig mode).
Stage values evaluate to:

\[
\begin{aligned}
k_1 &= z\,y, \\
k_2 &= z\,(y + (dt/2)\,k_1) = z\,y\,(1 + z\,dt/2), \\
k_3 &= z\,(y + (dt/2)\,k_2) = z\,y\,(1 + z\,dt/2 + (z\,dt)^2/4), \\
k_4 &= z\,(y + dt\,k_3) = z\,y\,(1 + z\,dt + (z\,dt)^2/2 + (z\,dt)^3/4).
\end{aligned}
\]

Substituting into the update:

\[
y^{n+1} \;=\; y^n\,\Bigl[1 + z\,dt + \tfrac{(z\,dt)^2}{2} + \tfrac{(z\,dt)^3}{6} + \tfrac{(z\,dt)^4}{24}\Bigr].
\]

This is exactly the \textbf{fourth-order Taylor truncation} of
\(e^{z\,dt}\). The next term, \((z\,dt)^5/120\), is the local truncation
error --- \passthrough{\lstinline!O(dt\^5)!} per step, hence
\passthrough{\lstinline!O(dt\^4)!} global error after
\passthrough{\lstinline!T/dt!} steps.

\subsubsection{Stability region intersection with the negative real
axis}\label{stability-region-intersection-with-the-negative-real-axis}

The amplification factor \(g(z\,dt)\) above must satisfy \(|g| \le 1\)
for stability. For real negative \(z\,dt = -x\) (\(x > 0\)),

\[
g(-x) \;=\; 1 - x + \tfrac{x^2}{2} - \tfrac{x^3}{6} + \tfrac{x^4}{24}.
\]

Setting \(g(-x) = -1\) and solving numerically yields
\(x \approx 2.7853\). For \(x \in [0,\,2.7853]\), \(g(-x) \in [-1,\,1]\)
and the method is stable; beyond that, \(|g| > 1\) and the discrete
solution amplifies.

For the heat equation with \(z = D\,\lambda_h\) and the worst-case
discrete eigenvalue \(\lambda_h^{\max} = -\sum_a 4/h_a^2\), the RK4
stability bound is

\[
4\,D\,dt\,\sum_a \frac{1}{h_a^2} \;\le\; 2.7853, \qquad
\boxed{\,D\,dt\,\sum_a \frac{1}{h_a^2} \;\le\; 0.6963.\,}
\]

This is the value encoded as
\passthrough{\lstinline!RK4::diffusionCFLLimit!}. Compared with forward
Euler's \passthrough{\lstinline!0.5!}, RK4 is
\passthrough{\lstinline!\~1.39×!} looser --- modest but useful when the
spatial grid is fine and Euler is at the edge of its CFL.

\subsubsection{Why combined refinement at CFL-limited dt is
spatial-dominated}\label{why-combined-refinement-at-cfl-limited-dt-is-spatial-dominated}

When \passthrough{\lstinline!dt = c · h\^2 / D!} (CFL-limited), the
time-error term scales differently for the two integrators:

\begin{itemize}
\tightlist
\item
  \textbf{Euler} time error per step: \(\sim |z|^2\,dt^2 / 2 = O(h^4)\)
  per step; over \(T/dt = O(h^{-2})\) steps, total time error
  \(O(h^2)\). Combined with spatial error \(O(h^2)\), total error is
  \textbf{\(O(h^2)\)}.
\item
  \textbf{RK4} time error per step: \(|z|^5\,dt^5 / 120 = O(h^{10})\)
  per step; over \(O(h^{-2})\) steps, total time error \(O(h^8)\).
  Spatial error remains \(O(h^2)\). Combined error is
  \textbf{\(O(h^2)\)} --- also spatial-dominated, but the constant is
  smaller.
\end{itemize}

So a combined-refinement convergence sweep (the user-chosen test shape)
shows slope \passthrough{\lstinline!\~2!} for \textbf{both} integrators;
the RK4 vs Euler benefit appears in the \textbf{constant}, not the
\textbf{slope}. The
\passthrough{\lstinline!IntegratorTest.RK4LessErrorThanEulerAtSameDt!}
test demonstrates that constant differs by at least 10× when the
time-error contribution is fully isolated from the spatial-error
contribution (using a pure-ODE manufactured solution with no Laplacian).

The roadmap-mandated \textbf{dt\^{}4 slope} is verified by
\passthrough{\lstinline!IntegratorTest.RK4TimeAccuracyOrder4!} --- also
a pure-ODE test, since the slope is hidden behind the spatial floor for
any PDE-based test under CFL-limited stepping.

\subsection{3. Algorithm}\label{algorithm}

\subsubsection{Tag types}\label{tag-types}

\href{../../../include/gridcalc/solve/Integrator.hpp}{\passthrough{\lstinline!include/gridcalc/solve/Integrator.hpp!}}
is an aggregator that pulls in
\href{../../../include/gridcalc/solve/ExplicitEuler.hpp}{\passthrough{\lstinline!ExplicitEuler.hpp!}}
and
\href{../../../include/gridcalc/solve/RK4.hpp}{\passthrough{\lstinline!RK4.hpp!}}.
Each per-integrator header defines its own tag struct:

\begin{lstlisting}[language={C++}]
struct ExplicitEuler { static constexpr double diffusionCFLLimit = 0.5;    };
struct RK4           { static constexpr double diffusionCFLLimit = 0.6963; };
\end{lstlisting}

This co-location avoids a circular include --- if the tags lived in
\passthrough{\lstinline!Integrator.hpp!} and
\passthrough{\lstinline!Integrator.hpp!} included the per-integrator
headers, then those per-integrator headers (which need the tag types)
couldn't include \passthrough{\lstinline!Integrator.hpp!} back without a
cycle.

\subsubsection{\texorpdfstring{Generic \texttt{integrate}
overloads}{Generic integrate overloads}}\label{generic-integrate-overloads}

Each per-integrator header defines its
\passthrough{\lstinline!integrate(...)!} overload tagged on its own tag
struct. So the dispatch is plain function-overload resolution at the
call site; no \passthrough{\lstinline!if constexpr!} ladder, no virtual
dispatch.

\passthrough{\lstinline!integrate(... RK4)!} (in
\passthrough{\lstinline!RK4.hpp!}) implements the four-stage update
using a small \passthrough{\lstinline!detail::axpyFresh(a, scale, b)!}
helper that returns \passthrough{\lstinline!a + scale * b!} element-wise
as a fresh \passthrough{\lstinline!Field!}. Per step, this allocates
four stage outputs (\passthrough{\lstinline!k1, k2, k3, k4!}) plus three
intermediate states (\passthrough{\lstinline!psi + (dt/2) k1!},
\passthrough{\lstinline!psi + (dt/2) k2!},
\passthrough{\lstinline!psi + dt k3!}) --- about six
\passthrough{\lstinline!Field<double>!} allocations per step at
\(N_x N_y N_z\) doubles each. Phase 20 will swap this for a persistent
scratch pool with a non-allocating RHS callback shape.

The final accumulation is fused into one element-wise loop over
\passthrough{\lstinline!(k1 + k4)/6 + (k2 + k3)/3!}, reducing reads from
the four \passthrough{\lstinline!k\_i!} arrays without an intermediate
sum field.

\subsubsection{\texorpdfstring{\texttt{explicitEuler} thin
wrapper}{explicitEuler thin wrapper}}\label{expliciteuler-thin-wrapper}

\href{../../../include/gridcalc/solve/ExplicitEuler.hpp}{\passthrough{\lstinline!ExplicitEuler.hpp!}}
keeps Phase 5's free function name:

\begin{lstlisting}[language={C++}]
template <class Rhs>
inline void explicitEuler(core::Field<double>& psi, Rhs&& rhs,
                          double dt, int n_steps) {
  integrate(psi, std::forward<Rhs>(rhs), dt, n_steps, ExplicitEuler{});
}
\end{lstlisting}

Phase 5 callers compile unchanged; the
\passthrough{\lstinline!DiffusionTest.GenericExplicitEulerOnZeroRhs!}
test continues to exercise the wrapper directly and
\passthrough{\lstinline!IntegratorTest.GenericIntegrateDispatchEulerEqualsExplicitEuler!}
verifies the wrapper produces bit-identical output to the underlying
\passthrough{\lstinline!integrate(..., ExplicitEuler\{\})!} call.

\subsubsection{\texorpdfstring{\texttt{diffuse} becomes a
template}{diffuse becomes a template}}\label{diffuse-becomes-a-template}

\href{../../../include/gridcalc/solve/Diffusion.hpp}{\passthrough{\lstinline!Diffusion.hpp!}}:

\begin{lstlisting}[language={C++}]
template <class Integrator = ExplicitEuler>
inline void diffuse(core::Field<double>& psi, double D, double dt,
                    int n_steps, Integrator tag = {}) {
  /* validate args */
  detail::checkDiffusionCFL<Integrator>(psi.getGrid(), D, dt);
  integrate(
      psi,
      [](const core::Field<double>& f) { return diff::laplacian(f); },
      D * dt,
      n_steps,
      tag);
}
\end{lstlisting}

The CFL check reads
\passthrough{\lstinline!Integrator::diffusionCFLLimit!} so each
integrator gets its own bound.
\passthrough{\lstinline!typeid(Tag).name()!} is included in the error
message for diagnostics.

\subsection{4. Design decisions}\label{design-decisions}

\begin{itemize}
\tightlist
\item
  \textbf{Tag dispatch over class hierarchy.} Compile-time dispatch
  matches the rest of the project (Phase 3's
  \passthrough{\lstinline!Pairwise!}/\passthrough{\lstinline!Kahan!},
  Phase 4's \passthrough{\lstinline!evaluate!} arity). No virtual-call
  overhead in the per-step inner loop. The tag's
  \passthrough{\lstinline!static constexpr!} properties (CFL limit) are
  available to higher-level drivers without any indirection.
\item
  \textbf{Tag types co-located with their integrators.} Avoids the
  circular include described in §3 above. The downside --- a caller
  writing \passthrough{\lstinline!solve::ExplicitEuler\{\}!} needs to
  have included \passthrough{\lstinline!solve/ExplicitEuler.hpp!} (or
  \passthrough{\lstinline!solve/Integrator.hpp!} which transitively
  pulls it) --- is acceptable. The \passthrough{\lstinline!diffuse!}
  driver pulls \passthrough{\lstinline!Integrator.hpp!} so any tag works
  through that entry point.
\item
  \textbf{\passthrough{\lstinline!explicitEuler!} kept as wrapper rather
  than removed.} Phase 5 callers and tests continue working unchanged.
  The two equivalent forms are documented as a stylistic choice.
\item
  \textbf{\passthrough{\lstinline!diffuse!} templated on integrator
  rather than two separate drivers.} One driver, one CFL-check
  implementation, parameterized on the tag. Adding a third integrator
  later (e.g., Phase 19's implicit BDF) needs only a new tag struct and
  a new \passthrough{\lstinline!integrate(... BDF)!} overload --- no
  change to \passthrough{\lstinline!diffuse!}.
\item
  \textbf{Pure-ODE convergence tests for \passthrough{\lstinline!dt\^4!}
  slope.} PDE-based tests at CFL-limited \passthrough{\lstinline!dt!}
  are spatial-dominated for both integrators (§2); the user's chosen
  combined-refinement test shows slope \passthrough{\lstinline!\~2!}
  (the expected combined order). To honor the roadmap's literal
  \passthrough{\lstinline!Δt\^4!} acceptance, an additional test uses a
  pure linear-decay ODE \passthrough{\lstinline!dpsi/dt = -α ψ!} with no
  spatial discretization, where the time-error slope is unambiguous.
\item
  \textbf{No persistent scratch pool yet.} Per-step allocation cost
  (\textasciitilde6 fields for RK4) is documented as a known cost. Phase
  20 (performance) is the right place to introduce a non-allocating RHS
  shape and a reusable scratch.
\end{itemize}

\subsection{5. References}\label{references}

{[}1{]} Hairer, E., Nørsett, S. P., \& Wanner, G. (2008). \emph{Solving
Ordinary Differential Equations I: Nonstiff Problems} (2nd revised ed.).
Springer Series in Computational Mathematics, vol.~8. Springer-Verlag.
ISBN 978-3-540-56670-0. \url{https://doi.org/10.1007/978-3-540-78862-1}.
Chapter II.1 covers the classical Runge--Kutta methods including the
specific 4-stage scheme used here, with the Butcher-tableau derivation
in §II.2 and the order-condition analysis in §II.3.

{[}2{]} Butcher, J. C. (2016). \emph{Numerical Methods for Ordinary
Differential Equations} (3rd ed.). Wiley. ISBN 978-1-119-12150-3.
\url{https://doi.org/10.1002/9781119121534}. The canonical modern
reference for Runge--Kutta theory; §3 gives the order-condition
framework that proves RK4's local truncation error is
\passthrough{\lstinline!O(dt\^5)!}.

{[}3{]} Wikipedia: \emph{Runge--Kutta methods} --- \emph{The
Runge--Kutta method} (RK4 section) and \emph{Stability} subsection.
\url{https://en.wikipedia.org/wiki/Runge\%E2\%80\%93Kutta_methods}
(permanent URL). Includes the explicit form of the RK4 update used in §3
and the stability-region plot showing the
\passthrough{\lstinline!\~−2.7853!} real-axis intersection.

{[}4{]} LeVeque, R. J. (2007). \emph{Finite Difference Methods for
Ordinary and Partial Differential Equations}. SIAM. ISBN
978-0-89871-783-9. \url{https://doi.org/10.1137/1.9780898717839}. §7.4
covers RK methods applied to method-of-lines discretizations of
parabolic PDEs, the framework used in
\passthrough{\lstinline!solve::diffuse!}. Cited from earlier phases as
well; this chapter is new context for Phase 6.
